% 13-operations.tex: surviving operational reality.
\section{Surviving reality: the stack, failure, load, the standing witness, and the day a pilot ends}
\label{sec:operations}

\motif{A correct system is not a system that survives.}

\analogy{A town that is right about everything and falls down in the first storm protects no one. So it rehearses its fires, keeps a second hall in another valley that takes no decisions until the first is gone, and measures how long its gates stay shut.}

\subsection{The stack}

The production topology is five services: a self-built TLS edge that negotiates a hybrid
post-quantum key exchange with modern clients, the application under gunicorn, a connection pooler,
PostgreSQL with continuous archiving, and Redis for the shared rate limiter. Every service runs as a
non-root user with every capability dropped, third-party images are pinned by digest, and every
production container drops the test framework. Four deployment paths share one secrets discipline,
one backup path and one threat model: a laptop launcher, a single-host compose profile with a
blue-green rolling deploy, a scripted Linux install under systemd exercised on two distributions,
and a Helm reference profile with default-deny network policies and the restricted pod security
standard, booted in CI on a one-node cluster. \sImpl

\subsection{Failure, drilled and measured}

The repository's rule is to exercise rather than read: ten of ten defects found in one sweep were
invisible to reading. Every operational claim is a drill that runs in CI and commits a number: high availability under
a live write stream, the loss of a region, replicas, partitions, backup and restore, chaos under
traffic and rolling deploys, each against a ceiling. The drills and their numbers are listed in
Appendix~\ref{app:status}.


\subsection{Load: what was measured, what was projected, and what a number rests on}

Three kinds of number appear in the repository, and the paper keeps them apart. \emph{Measured}
numbers come from a named machine, an Apple Silicon laptop with eight cores, and carry a version
stamp. \emph{Simulated} numbers come from the national simulation harness, a synthetic country
driven through the real procedures at a stated scale factor. \emph{Projected} numbers are arithmetic
over measured ones and are labelled as such by the repository itself, under a named assumption. No
production cluster has been measured.

Table~\ref{tab:perf}, in Appendix~\ref{app:status}, lists every performance number with its kind and
the version that measured it.


The roadmap's planning targets (ten million persons for a state; 350 million, five thousand
sustained and fifty thousand peak verifications per second nationally) were checked
against a computational throughput model, and both halves of the answer are recorded.
The model is built from single-machine measurements under stated assumptions; nothing has run at
national scale, and the first bottleneck no measurement reaches is operational (replication,
pooling and partition maintenance under a real multi-instance load). Every throughput target clears by more
than an order of magnitude: the peak needs about six and a half cores of single-witness
verification, the sustained rate under one, and a surge of two hundred thousand enrolments a day is
nine minutes of one signer. And the system could not have run for a week at the sustained rate,
because the verification event's identifier was a 32-bit sequence: two billion rows at five thousand
a second is five days, twelve hours at the peak, and nothing is slow when a sequence is exhausted,
every insert on the path simply fails. Five surrogate identifiers are 64-bit now, a check recomputes
the arithmetic from the live schema on every push, and the rule the repository drew from it is that a
capacity question answered in cores and seconds is the easy half. The model also refuses to call a
core count \emph{measured}: the per-core rate is measured, the multiplication has never been run, and
every extrapolation carries its assumption by name. A separate cost model
finds that verification throughput is not the cost driver at any realistic national scale, since a
hundred million people verified twelve times a year is thirty-eight verifications per second on
average, and states which of its inputs are facts about the code and which are guesses about a
deployment. A national workload has been simulated at a scale factor and driven through the real
constraints; it has not been run. \sExp

\subsection{Abuse, secrets, supply chain, and decisions that leave a record}

Per-agency quotas on issuance, revocation and verification are opt-in rows enforced by a trigger on
every write path, exact under concurrency by an advisory lock, refused loudly on every surface, and
keyed by agency identifiers, never by a person; a quota is a decision with a
reason that survives being changed rather than a setting that can be overwritten. Velocity alerts
compare each agency to its own trailing week. The per-IP rate limiter shares one Redis bucket across
workers. Operators authenticate with WebAuthn as a second factor after an enrolment deadline,
sessions are registered server-side with per-role caps and origin checks, and a per-role network
policy applies. Creating an operator account leaves a record, and every audit event type the schema names has a
writer, pinned by a check. Secrets are file-mounted, sealed and rotated through one
lifecycle, and a database password reaches the database before its file changes, or is not rotated at
all; the application refuses to boot in production on a default key, and the application log
records no form body, cookie value or token value, which a check holds by reading the logging
calls. Every release ships bills of materials with
signed provenance, and CVE scans of the dependency surface and of every self-built image gate the
build. \sImpl

\subsection{The standing witness, and whether it watches}

Twenty CI jobs run on every push, and most carry a comment naming the past failure they exist to
prevent. They run the product suite and the invariant layer; drive the Atlas console in a headless
browser; build and boot the images; build the edge and prove the post-quantum handshake against a
real certificate; page a duress alert through a real alert manager to a webhook; validate the
tracing wire and the dashboards; kill and restore the primary; install the Linux profile on two
distributions; sign in-token through a software PKCS\#11 module with two witnesses; roll a deploy
under traffic; evacuate a region; fail the HA profile's leader, lease store and etcd member under a
live write stream; boot the Kubernetes profile and fail its database over; sign and verify with real
ML-DSA-65; federate two instances over HTTP; install \code{polaris-verify} as a stranger would; type-check,
test and conform the TypeScript SDK; scan dependencies and every image for vulnerabilities; and boot
the full production stack end to end. \sImpl

The drills are themselves watched. A drill refuses to report a guarantee it did not test, and a check
holds every drill to counting its cases; every CI job installs what it runs, so a job fails on a
finding rather than on its own setup; and the local gate a maintainer runs before pushing runs what
CI runs, names the verification a change needs from the paths and route handlers that moved, and
repeats the web suite as the production role and under the real signer. \sImpl

\subsection{The day a pilot ends, and the day the old credential stops being accepted}

Two operations were written down as procedures that run rather than as paragraphs, because each is
the promise an institution actually says yes on. A pilot's real promise is not that it will work but
that it can be wound back, and that promise fails quietly: erasure becomes a paragraph in a consent
form nobody executes. The wind-down revokes every pilot credential and replaces every
name with a meaningless marker under the append-only erasure record, the consent language a
participant is shown is generated from what the code does rather than written from what would be
reassuring, and the whole procedure runs against a real database on every push. The pilot pack that
an institution owes before it starts is closed with it. \sImplMark{} for the procedure; \sExt{} for a
pilot, since none has run.

A new national credential arrives beside a driving licence and a passport, and for years it is the
second thing somebody carries. The dangerous moment in that period is not the cutover but the
sunset, the day the old credential stops being accepted, because that is the moment the new one
becomes compulsory without anybody having decided to make it so. The coexistence plan
treats the sunset as a decision with a floor under it rather than a phase that arrives, and its
verdict refuses to be computed from what the issuer can see: whether relying parties accept the
credential and whether an alternate path exists and is usable are attestations an operator makes,
and the system records them rather than guessing. \sImpl

Two mappings sit beside these. The NIST SP 800-63 mapping states, row by row, what an assurance
claim would rest on, every citation resolved and run by a drill so that a row cannot claim evidence
that is not there; it is explicitly not a conformance claim and no assessment has been performed.
And accessibility is enforced on every push for the third that automation covers, and named as not
done for the rest, because an identity system a person cannot operate is one that excludes them from
identity. \sImplMark{} for the mechanisms; \sExt{} for an assessment.

\begin{layercard}{Layer card: operations}
\qa{What it is}{The five-service stack on four substrates; HA, a standby region, replicas, partitioning, backup and disaster recovery, chaos, rolling deploys; abuse controls, secrets, supply-chain gates; twenty CI jobs and two scheduled drills; the pilot wind-down and the coexistence plan.}
\qa{Why it exists}{A correct system that cannot survive a lost disk, a lost region, a bad deploy or a hostile client is not infrastructure, and one that cannot be wound back is not one an institution should start.}
\qa{What it trusts}{The drills' ephemeral hosts as a lower bound on production behaviour; the operator's placement of hosts and regions and the custody choice.}
\qa{What it refuses}{A claim without a drill; a number without a stamp; a core count called measured; a drill that reports a guarantee it did not test; a root-running container; an unpinned image; a fixable critical CVE.}
\qa{What it can see}{Metrics, traces and alerts about the system. The correlation id on every request never touches the audit of record.}
\qa{What it cannot see}{A person: quotas, alerts, metrics and logs are keyed by agency, route and worker.}
\qa{If it fails}{A drill that misses its ceiling fails the build; a chaos or DR finding becomes a check; a region evacuation that invents or skips a row fails; a production incident is what the readiness ledger's nine operator decisions exist to prepare for.}
\qa{Now or later}{Now, on drills. A hardware module, production hardware and a real pilot are later.}
\qa{Inside and outside}{Inside: everything above. Outside: the system's view of itself, which an operator needs without a view of people.}
\end{layercard}

\nextlayer{An operator running all of this needs to see what the system is doing and to understand
where authority comes from, and must be able to do both without the system becoming a map of its
people. The next layer is how Polaris observes and understands itself.}
